\chapter{Fundamentals and Related Work}
\label{cha:Fundamentals}

Before starting with the introduction to the InpaintAR tool some general knowledge and information needs to be further elaborated upon.
All of the concepts relevant for the tool, and which algorithms and logics it uses to achieve its goal of reducing the users distraction are being presented, starting from a more detailed explanation on Visual Clutter and why it is relevant in MR, 
to a more detailed introduction to DR, inpainting, how those two concepts can be combined and the inpainting methods analyzed in the tool and this thesis.

\section{What is Visual Clutter}
Clutter can be seen as a state in which items, their placement or some features causes performance degradations in certain tasks. 
It can lead to a decrease in object recognition performance or cause issues when searching for certain objects. In addition clutter also seems to have
a detrimental effect to short-term memory, whereas not only the number of elements, but their features can have a negative effect aswell \cite{RosenholtzRuth2007Mvc}.

\begin{figure}[htb!]
  \begin{center}
    \includegraphics[width=0.45\textwidth]{ClutterToReactionTime.PNG}
    \caption{Taken from \cite{RosenholtzRuth2007Mvc}, Page 3; It shows a comparison between reaction time and clutter to predict performance. The difficulty of the given task may also change the increase in reaction time compared to the clutter, therefore 3 curves are provided instead of one.}
    \label{fig:ClutterToReactionTime}
  \end{center}
\end{figure}

\subsection{Why does it matter for Mixed Reality}
Mixed Reality (MR), as defined by Paul Milgram in 1994, can be seen as an environment in which real and virtual objects are presented together through a single display. Therefore AR, DR and Virtual Reality (VR) can all be seen as subcategories of MR, which are placed somewhere in the virtuality continuum which spans a completely real environment to a completely virtual one \cite{Milgram1994MR}.

As defined by S. Mann and J.Fung back in 2002, AR aims to enhance reality through the addition of virtual objects, whilst DR aims to remove physical objects from the users perception \cite{Mann2002}.

Whilst during the design of traditional software and information visualization, information overload and clear user interfaces need to be carefully considered.
In these domains all of the factors which may induce clutter and information overload can be controlled internally in the visualization or software \cite{RosenholtzRuth2007Mvc}.

AR is uniquely flawed in that regard due to it enhancing the surrounding environment of the user. This in turn means there could be a lot of external factors which may cause clutter, starting from the pure concept and interactivity of AR
oftentimes allowing objects to be placed freely in the environment and therefore cluttered environments to the space constraints of the user themselves.

\begin{figure}[htb!]
  \begin{center}
    \includegraphics[width=0.5\textwidth]{clutteredChess.png}
    \caption{\textbf{TODO SG; Add after UseCase Example done;} An example of virtual objects enhancing an already cluttered environment and therefore worsening the existing visual clutter}
    \label{fig:clutteredChess}
  \end{center}
\end{figure}

\section{Diminished Reality}
qwerqwerqwerqwer

\section{Inpainting}

\subsection{What is inpainting}

\subsection{Methods analyzed in this thesis}

\subsubsection{Fast Marching Method}

\subsubsection{Texture Synthesis and Region Segmentation}

\subsubsection{Nonlocal Texture Matching and nonlinear filtering}
Still TBD, depends if too similar to Texture Synthesis or not

\subsubsection{Nearest Neighbor Pixel Filling}

\subsubsection{DeepLearning}

\subsubsection{Regression}

\subsection{Inpainting in Diminished Reality}
